\section*{Comparative Reception and Theological Dossier}

\begin{studybox}{The question beneath the packet}
Year C does not ask whether work, property, planning, or enjoyment are evil.
It asks what they can bear. Qoheleth tests labor against succession and death;
the psalm asks for wisdom and established work; Luke tests abundance against
the same night; Colossians locates life in the risen Christ. Reception is
useful only when it preserves those distinctions.
\end{studybox}

\subsection*{\emph{Hevel}: not nothingness, but the failure of mastery}

The first reading joins Ecclesiastes~1:2 to 2:21--23. That discontinuity creates
a thesis-and-case relation: the superlative \emph{hevel} saying is followed by
a worker who has used wisdom, knowledge, and skill yet must leave the result
to someone who did not labor for it. Anxiety reaches into the night. The
problem is not simply that the successor might be foolish; the worker cannot
make achievement a permanent possession or control what happens after death.

“Vanity” has a long Christian afterlife, but it can mislead if heard only as
moral disgust or ontological worthlessness. Breath, vapor, transience, and
enigma better expose the limit at issue. Qoheleth elsewhere recognizes
received enjoyment as gift; this cut therefore should not be made to condemn
work or created pleasure absolutely. It condemns their conscription as
guarantees against mortality.

The discontinuous lection also creates a risk: the refrain can swallow the
specific labor argument. A responsible reading keeps both. Everything cannot
be mastered, and here the concrete proof is work transferred beyond the
worker's control.

\subsection*{Psalm 90: numbering days without despising work}

Psalm~90 answers neither with accumulation nor abandonment. Human beings return
to dust; generations pass before God like a watch in the night; life resembles
grass. Yet the prayer asks for a wise heart, returning mercy, morning
satisfaction, gladness proportioned to affliction, divine favor, and twice for
the work of human hands to be established.

Augustine receives the psalm within Christian hope, reading mortality,
conversion, mercy, and established work through an ecclesial and
Christological horizon. His numerical constructions are historical reception,
not definitions of the psalm's Hebrew. The most important liturgical effect is
simpler: mortality does not make work worthless. It makes wisdom, mercy, and
divine establishment necessary.

The official occurrence page's response-marker discrepancy must remain
visible. The response words come from Psalm~95:8, not Psalm~90:1. Its
imperative “today” intensifies the numbered-days petition, but the guide does
not silently correct a protected book it has not collated at typography level.

\subsection*{The rich fool: grammar, audience, and social absence}

Luke distinguishes an inheritance petitioner, Jesus' answer to him, the crowd
addressed by the warning, and the parable's landowner. Jesus neither decides
nor narrates the merits of the brothers' civil dispute. The crowd receives the
warning because greed exceeds one lawsuit. The parable then makes abundance
grammatical: the man speaks to himself, asks what he will do, and repeats
``my'' barns, grain, goods, and soul. Land has produced plentifully, but giver,
worker, neighbor, poor person, and heir are absent from his deliberation.

God's address reverses the imagined timetable. “Many years” becomes “this
night”; the soul commanded to rest becomes a life demanded; the goods
possessed become goods transferred to an unnamed other. The parable does not
condemn barns as buildings or prudence as a virtue. It exposes a plan closed
to relation, accountability, mortality, and God.

Augustine's sermon preserves the movement from inheritance to greed and death
but adds a judgment about the petitioner's civil case that Luke does not
supply. Ambrose relocates the barns into a famine and grain-withholding
argument. Basil's homily makes surplus morally answerable to the hungry and
unclothed. These social readings disclose a genuine implication of
self-enclosed abundance, but their reconstructed settings must not be smuggled
back into Luke's unnamed harvest.

Francis's 2019 Angelus offers a modern pastoral reception with unusually close
attention to the text's sequence: inheritance quarrels, accumulated goods,
imagined years, the interruption of “this night,” enslavement to possessions,
and sharing. He explicitly says material goods are goods and necessary means,
while denying them the capacity to ground life.

\subsection*{Colossians: greed as idolatry and the renewed social person}

Colossians is not the product of official Old Testament--Gospel correlation.
Its semi-continuous course nevertheless creates a powerful canonical overlay.
Those raised with Christ seek things above because their life is hidden with
him and will appear with him. The lection then names greed as idolatry and
moves, after omitted verses, to stripping off the old person, renewal in the
creator's image, and a humanity in which Christ is all and in all.

Chrysostom refuses a spatial escapism: “above” directs life toward Christ and
therefore changes conduct. The immediate omitted verses name anger, malice,
slander, and abusive speech; their absence from proclamation cannot make them
irrelevant to interpretation. Greed is idolatry not because matter is evil but
because desire grants a created good the allegiance and security owed to God.

The renewed person is also social. Greek and Jew, circumcised and
uncircumcised, barbarian, Scythian, slave, and free cannot serve as rival
grounds of worth in the new humanity. This answers the rich man's isolated
possessive grammar more deeply than a private counsel to feel detached.

\subsection*{Poverty of spirit, sharing, and the danger of romanticizing need}

The Gospel acclamation gives an authorized liturgical lens: poverty of spirit
and the kingdom. It cannot be reduced to involuntary material deprivation or
used to bless conditions that deny food, housing, wages, or protection.
Neither may “spiritual” poverty be used to erase the Gospel's material
warning. The beatitude names humble receptivity before God; Luke tests whether
actual goods are enclosed within the self.

The two Communion alternatives sharpen this relation differently. Wisdom's
heavenly food makes nourishment gift; John's bread of life makes the giver
personally central. Against the barns, either ending makes food receptive and
relational. Yet the alternatives remain alternatives, and neither proves an
economic policy from the Missal.

\subsection*{Changed-register reception: image and music}

Rembrandt's 1627 \emph{Parable of the Rich Man} abandons field and barn for a
nocturnal interior where a solitary figure examines coins by artificial
light. The changed setting interprets the parable through bookkeeping,
attention, and isolation. It is an artistic afterlife, not evidence about
first-century storage.

Charles Ives's \emph{Psalm 90} carries the psalm into modern concert and
choral practice, juxtaposing human frailty, divine eternity, and the
establishment of work through musical recurrence and massed sound. Again the
change of register matters: a score can make temporal contrast audible but
cannot adjudicate the psalm's authorship or Hebrew semantics.

The 2014 \emph{Homiletic Directory} supplies an official reception of another
kind. It identifies the rich fool within Luke's Year C mercy-and-warning
course, explains why Old Testament and Gospel are correlated, and preserves
the apostolic reading's independent semi-continuous path. That structural
discipline is itself a correction to homilies that force every reading into a
single historical theme.

\begin{threestable}{Tension}{Responsible synthesis}{Misuse excluded}
Work and death & Work is real good received under mortality and divine judgment. & Nihilism, idleness, or treating productivity as immortality.\\
Property and gift & Goods can serve honest life, dependents, solidarity, and the poor. & Absolute condemnation of ownership or self-enclosed accumulation.\\
Hidden life and public duty & Union with Christ renews desire, speech, and social belonging. & Escapism or a merely private “spiritual” detachment.\\
Poverty and Eucharistic reception & The creature receives life and food from a giver and is turned toward communion. & Romanticizing involuntary poverty or collapsing sacrament into economics.\\
\end{threestable}
